\section{Figures}

\subsection{Figure Between Paragraphs}

Paragraph before the figure.

\begin{figure}
  \includegraphics[width=\linewidth, height=3.5cm]{example-image-a}
  \caption{Single full-width figure between two paragraphs.}
  \label{fig:t-single}
\end{figure}

Paragraph after the figure, referring to \cref{fig:t-single}.

\subsection{Figure Directly After a Heading}

\begin{figure}
  \includegraphics[width=0.6\linewidth, height=3cm]{example-image-b}
  \caption{A 60\% figure right after a heading.}
\end{figure}

\subsection{Narrow Figures: Widths and Captions}

\begin{figure}
  \includegraphics[width=0.3\linewidth]{example-image-a}
  \caption{30\% width, short caption.}
\end{figure}

\begin{figure}
  \includegraphics[width=0.5\linewidth]{example-image-b}
  \caption{50\% width with a long caption that does not fit on one line, so it
           is justified across the full text width (IEEE style).}
\end{figure}

\begin{figure}
  \includegraphics[width=0.75\linewidth, height=4cm]{example-image-c}
  \caption{75\% width.}
\end{figure}

\begin{figure}
  \includegraphics[height=6cm]{example-image-9x16}
  \caption{Tall portrait image (sized by height), with a caption long enough to
           need two lines below a narrow image.}
\end{figure}

Text after the narrow figures.

\subsection{Full Width: Short and Long Captions}

\begin{figure}
  \includegraphics[width=\linewidth, height=3cm]{example-image-a}
  \caption{Short caption under a full-width figure.}
\end{figure}

\begin{figure}
  \includegraphics[width=\linewidth, height=3cm]{example-image-b}
  \caption{A long caption under a full-width figure. It explains what the
           colours mean and what the reader should notice, so it wraps onto a
           second line.}
\end{figure}

\subsection{Two Figures in a Row}

\begin{figure}
  \includegraphics[width=\linewidth, height=2.5cm]{example-image-a}
  \caption{First of two figures.}
\end{figure}

\begin{figure}
  \includegraphics[width=\linewidth, height=2.5cm]{example-image-b}
  \caption{Second figure directly after the first.}
\end{figure}

\subsection{Side by Side}

\begin{figure}
  \begin{subfigure}[t]{0.48\linewidth}
    \includegraphics[width=\linewidth, height=3cm]{example-image-a}
    \caption{Left image.}
    \label{fig:t-left}
  \end{subfigure}\hfill
  \begin{subfigure}[t]{0.48\linewidth}
    \includegraphics[width=\linewidth, height=3cm]{example-image-b}
    \caption{Right image, with a caption long enough to wrap onto two lines.}
    \label{fig:t-right}
  \end{subfigure}
\end{figure}

\Cref{fig:t-left,fig:t-right} are side by side.

\subsection{Grid of Four}

\begin{figure}
  \begin{subfigure}[t]{0.48\linewidth}
    \includegraphics[width=\linewidth, height=2.5cm]{example-image-a}
    \caption{Top left.}
  \end{subfigure}\hfill
  \begin{subfigure}[t]{0.48\linewidth}
    \includegraphics[width=\linewidth, height=2.5cm]{example-image-b}
    \caption{Top right.}
  \end{subfigure}
  \subfiguregap
  \begin{subfigure}[t]{0.48\linewidth}
    \includegraphics[width=\linewidth, height=2.5cm]{example-image-c}
    \caption{Bottom left.}
  \end{subfigure}\hfill
  \begin{subfigure}[t]{0.48\linewidth}
    \includegraphics[width=\linewidth, height=2.5cm]{example-image}
    \caption{Bottom right.}
  \end{subfigure}
\end{figure}

\subsection{Sub-figures With a Main Caption}

\begin{figure}
  \begin{subfigure}[t]{0.48\linewidth}
    \includegraphics[width=\linewidth, height=2.5cm]{example-image-a}
    \caption{Day.}
  \end{subfigure}\hfill
  \begin{subfigure}[t]{0.48\linewidth}
    \includegraphics[width=\linewidth, height=2.5cm]{example-image-b}
    \caption{Night.}
  \end{subfigure}
  \caption{A main caption for the whole figure, below its sub-figures.}
\end{figure}
